\chapter{Chapter 9 Refresher: Alignment, Instruction Tuning, and Preference Optimization --- Preferences, KL Regularization, and Policy Updates}

\scenelabel
Alignment reframes Preface Steps~2--4: instead of conditioning only on observed tokens, we condition on human comparisons and regularise policies toward trusted references.

\paragraph{Step A --- Pairwise likelihood.}
\textbf{Problem.} Scalar correctness insufficient; humans supply partial orders.
\textbf{Insight.} Bradley--Terry models lift score gaps into Bernoulli probabilities---Preface Step~2 specialised to preference events.
\textbf{Lineage.} Thurstone/B Bradley--Terry tradition; Rafailov et al.\ DPO (2023) streamlines derivatives.
\textbf{Mental model.} Relative labels estimate latent utilities; absolute calibration remains optional extra.

\paragraph{Step B --- KL anchor.}
\textbf{Problem.} Unconstrained policy optimisation hallucinates unsupported modes.
\textbf{Insight.} KL penalties replay Preface Step~10 mismatch narrative between \(\pi_\theta\) and \(\pi_{\mathrm{ref}}\).
\textbf{Lineage.} Schulman/PPO lineage intersects RLHF practice (Stiennon, Ouyang, subsequent stacks).
\textbf{Mental model.} \(\beta\) trades helpfulness vs behavioural drift---no universal optimum without governance context.

\paragraph{Step C --- Support hazards.}
\textbf{Problem.} Zero reference probability yields exploding log-ratios.
\textbf{Insight.} Conditioning events need positive measure---same vigilance as Preface Step~3 warned for chain factors.

Preference optimisation is probabilistic inference under sociotechnical data---mathematics unchanged, likelihood rewritten.

\section*{Notation Introduced in This Chapter}
\begin{itemize}[leftmargin=1.5em]
  \item $x$: prompt/context.
  \item $y^+, y^-$: preferred and dispreferred responses.
  \item $r(x,y)$: reward/score function.
  \item $\sigma(\cdot)$: logistic sigmoid function.
  \item $\pi$: current policy distribution.
  \item $\pi_{\mathrm{ref}}$: reference policy.
  \item $\beta$: KL regularization coefficient.
\end{itemize}

\section*{What You Need To Recall Precisely}
\begin{itemize}[leftmargin=1.5em]
  \item \textbf{Bradley--Terry form:}
  \[
    P(y^+ \succ y^- \mid x)=\sigma(r(x,y^+)-r(x,y^-)).
  \]
  Write margin \(\Delta r = r(x,y^+)-r(x,y^-)\). In Bradley--Terry style models this acts like a log-odds difference: \(\sigma(\Delta r)\) matches how humans choose stochastically under utility noise; pairwise losses therefore encourage reward gaps consistent with observed preferences.

  \item \textbf{KL regularization role.}
  Optimising preferences alone can peel \(\pi_\theta\) away from pretrained \(\pi_{\mathrm{ref}}\) into unsupported linguistic regimes where rewards mis-generalise. Adding \(\beta\, D_{\mathrm{KL}}(\pi_\theta \,\|\, \pi_{\mathrm{ref}})\) (or sampled surrogates thereof) penalises drift measured against data-supported behaviour while pairwise margins tilt \(\pi_\theta\) toward helpful completions.

  \item \textbf{Pairwise objective geometry:} score differences (margins) drive learning signal.
  \item \textbf{Support condition:} near-zero reference probability destabilizes log-ratio terms.
  \item \textbf{Objective-vs-behavior distinction:} optimizing preference loss does not automatically guarantee safety or honesty on all distributions.
\end{itemize}

\section*{How These Results Build on Each Other}
\begin{enumerate}[leftmargin=1.5em]
  \item Convert pairwise labels into likelihood terms.
  \item Learn scoring or policy update from pairwise margins.
  \item Add KL regularization to bound policy drift.
  \item Validate on held-out safety/helpfulness metrics.
\end{enumerate}

\section*{Reminder Glossary (Term by Term)}
\begin{itemize}[leftmargin=1.5em]
  \item \textbf{Preference pair:} chosen vs rejected response for same prompt.
  \item \textbf{Margin:} score difference between preferred and non-preferred outputs.
  \item \textbf{KL penalty:} regularizer limiting divergence from reference policy.
  \item \textbf{Support mismatch:} missing probability mass causing unstable ratios.
  \item \textbf{Over-optimization:} objective improves while true quality degrades.
\end{itemize}

\section*{Worked Mini Example (Pairwise Probability)}
If $r(x,y^+)=2.0$ and $r(x,y^-)=0.5$:
\[
P(y^+\succ y^-|x)=\sigma(1.5)\approx0.817.
\]
Interpretation: the model assigns roughly $82\%$ preference probability to $y^+$ under this reward gap.

\section*{Worked Example (Preference Gain vs Safety Drift)}
Suppose after tuning:
\begin{itemize}[leftmargin=1.5em]
  \item preference win-rate improves from 62\% to 73\%,
  \item refusal-on-unsafe-prompts worsens from 91\% to 84\%.
\end{itemize}
Reasoning: preference objective improved but safety behavior regressed.  
Implication: keep preference metric and safety metric as separate acceptance gates; one does not substitute for the other.

\section*{Intuition Checkpoints}
\begin{itemize}[leftmargin=1.5em]
  \item Lower KL does not mean better helpfulness; it means less policy drift.
  \item Better reward-model fit does not guarantee lower downstream harm.
  \item Pairwise preferences encode relative order, not absolute truth.
\end{itemize}

\section*{Further Refresh}
\begin{itemize}[leftmargin=1.5em]
  \item KL/cross-entropy foundations: see Chapter 1 refresher.
  \item Optimization mechanics: see Chapter 2 refresher.
  \item Adaptation substrate (LoRA/QLoRA): see Chapter 8 refresher.
  \item \textbf{Christiano/Stiennon/Ouyang papers} --- RLHF workflow and assumptions.
  \item \textbf{Rafailov et al. (DPO)} --- direct preference objective derivation.
\end{itemize}
